\chapter{复变函数基础}

\section{实数域到复数域的扩充}
从实数域 $\mathbb{R}$ 到复数域 $\mathbb{C}$ 的扩充，是为了使代数方程总有解、使开方、指数、三角函数等运算在更大范围内封闭。

\subsection{扩充动机}
实数域中，方程
$$x^2 + 1 = 0$$
无解。为使这类方程有解，引入虚数单位 $i$，满足
$$i^2 = -1.$$

\subsection{扩充步骤}
\begin{tcolorbox}[title=复数的定义]
	1. 引入符号 $i$，定义 $i^2 = -1$。
	
	2. 定义复数为 $z = x + iy$，其中 $x,y\in\mathbb{R}$。
	
	3. 规定复数的加法与乘法运算满足交换律、结合律、分配律。
	
	4. 当 $y=0$ 时，$z = x$，即实数是复数的特例。
	
	5. 当 $x=0$ 时，$z = iy$ 为纯虚数。
\end{tcolorbox}

\subsection{扩充后的封闭性}
\begin{tcolorbox}[title=复数域的封闭性, colframe=gray!60!black, colback=gray!5!white]
	扩充后的复数域 $\mathbb{C}$ 具有以下重要性质：
	\begin{enumerate}
		\item 复数域是\textbf{代数闭域}：任意次数的复系数多项式方程在 $\mathbb{C}$ 内必有根。
		\item 加、减、乘、除（除数非零）运算封闭。
		\item 开方、指数、对数、三角函数、反三角函数均可合理定义。
	\end{enumerate}
\end{tcolorbox}

\section{定义与性质}
\begin{tcolorbox}[title=复变函数定义, colframe=blue!50!black, colback=blue!5!white]
	复变函数是定义在复数集上的函数，记为
	$$w = f(z),$$
	其中 $z\in\mathbb{C}$ 为复自变量，$w\in\mathbb{C}$ 为复因变量。
	
	将 $z=x+iy$，$w=u+iv$ 代入，则复变函数可表示为
	$$f(z) = u(x,y) + i\,v(x,y),$$
	即一个复变函数等价于\textbf{两个二元实变函数的有序组合}。
\end{tcolorbox}

\section{初等函数}
\subsection{指数函数}
\begin{tcolorbox}[sharp corners]
	复指数函数定义为
	$$e^z = e^{x+iy} = e^x(\cos y + i\sin y),$$
	满足：
	\begin{itemize}
		\item 周期性：周期为 $2\pi i$；
		\item $e^{z_1+z_2}=e^{z_1}e^{z_2}$；
		\item $|e^z|=e^x>0$。
	\end{itemize}
\end{tcolorbox}

\subsection{对数函数}
\begin{tcolorbox}[sharp corners]
	复对数函数是指数函数的反函数，定义为多值函数：
	$$\ln z = \ln|z| + i\left(\arg z + 2k\pi\right),\quad k\in\mathbb{Z}.$$
	主值分支为
	$$\mathrm{Ln}\,z = \ln|z| + i\mathrm{Arg}\,z.$$
\end{tcolorbox}

\subsection{三角函数}
\begin{tcolorbox}[sharp corners]
	由欧拉公式定义复三角函数：
	$$\cos z = \frac{e^{iz} + e^{-iz}}{2},\quad \sin z = \frac{e^{iz} - e^{-iz}}{2i}.$$
	性质：
	\begin{itemize}
		\item 周期仍为 $2\pi$；
		\item 在复平面上\textbf{无界}；
		\item 实三角函数公式全部形式保留。
	\end{itemize}
\end{tcolorbox}

\subsection{反三角函数}
\begin{tcolorbox}[sharp corners]
	以反正弦函数为例，由 $\sin w = z$ 解出：
	$$\arcsin z = -i \ln\left(iz + \sqrt{1-z^2}\right).$$
\end{tcolorbox}

\subsection{一般幂函数 $z^b$}
\begin{tcolorbox}[sharp corners]
	利用复对数统一定义：
	$$z^b = e^{b\ln z}.$$
\end{tcolorbox}